\chapter{Linear Time Series}\label{ch:qm:linear-time-series}

The spread between the ten-year and the two-year Treasury yields, the 2s10s that a steepener trades, has a daily autocorrelation of
0.9987 over fifty years. A regression of tomorrow's spread on today's gives a slope of 0.99875 with a \omterm{def:qm:estimation:estimator}{standard error} of 0.00045: a
$t$-statistic of $-2.79$ against a slope of one, a one-sided $p$-value of 0.003 by the normal table. The spread seems to mean-revert,
with a \omterm{def:qm:stochastic-differential-equations:ou}{half-life} of 552 trading days. But if the spread had a \omterm{def:qm:linear-time-series:unitroot}{unit root}, the normal table would be the wrong one: the right one puts the
5\% critical value at $-2.86$, and the evidence falls short. This chapter is about the models that describe such series and the tests that
tell them apart: stationarity and the \omterm{def:qm:linear-time-series:acf}{autocorrelation function}, autoregressive and moving-average models, the spectral view, \omterm{def:qm:linear-time-series:unitroot}{unit roots}
and the regressions they make spurious, and \omterm{def:qm:linear-time-series:longmem}{long memory}.

\section{Stationarity and autocorrelation}

\begin{definition}[Stationary process, weak stationarity, white noise]\label{def:qm:linear-time-series:stationary}
A process $(X_t)$ is a \emph{stationary process}\index{stationary process} (strictly) if the joint law of $(X_{t_1 + h}, \dots, X_{t_k + h})$ does not
depend on $h$. It has \emph{weak stationarity}\index{weak stationarity} if its mean is constant and $\Cov(X_t, X_{t+h})$ depends only on $h$.
\emph{White noise}\index{white noise} is a weakly stationary sequence with mean zero and no autocorrelation at any nonzero lag.
\end{definition}

\begin{definition}[Autocovariance, autocorrelation and partial autocorrelation functions]\label{def:qm:linear-time-series:acf}
For a weakly \omterm{def:qm:linear-time-series:stationary}{stationary process}, the \emph{autocovariance function}\index{autocovariance function} is $\gamma(h) = \Cov(X_t, X_{t+h})$, and the
\emph{autocorrelation function}\index{autocorrelation function} is $\rho(h) = \gamma(h)/\gamma(0)$. The \emph{partial autocorrelation
function}\index{partial autocorrelation function} at lag $h$ is the last coefficient $\phi_{hh}$ of the best linear predictor of $X_{t+h}$ from $X_{t+h-1},
\dots, X_t$: the correlation at lag $h$ once the intermediate lags are accounted for.
\end{definition}

Under \omterm{def:qm:linear-time-series:stationary}{white noise} the sample autocorrelations are approximately independent $\mathcal N(0, 1/n)$, whence the $\pm 2/\sqrt n$ bands of every
correlogram. The decomposition that justifies the models of the next section is Wold's (1938): every weakly stationary, purely
nondeterministic process is an infinite moving average $X_t = \mu + \sum_{j \ge 0}\psi_j\varepsilon_{t-j}$ of its one-step prediction errors
$\varepsilon_t$, a \omterm{def:qm:linear-time-series:stationary}{white noise}, with $\psi_0 = 1$ and $\sum_j\psi_j^2 < \infty$.

The 2s10s spread's correlogram (\cref{fig:qm:linear-time-series:acf}) is the signature of a nearly nonstationary level: autocorrelations
that start at 0.9987 and fall almost linearly over a year of lags. Its daily changes are nearly \omterm{def:qm:linear-time-series:stationary}{white noise}, with a standard deviation of 4.6
basis points and autocorrelations of 0.015, 0.015 and $-0.027$ at the first three lags, within the band $\pm 0.018$ for the first two.

\begin{omfigure}
\begin{tikzpicture}
\begin{groupplot}[group style={group size=2 by 1,horizontal sep=1.5cm},width=7.4cm,height=4.8cm,
  tick label style={font=\small},label style={font=\small},grid=major,grid style={gray!25},table/col sep=comma]
\nextgroupplot[xlabel={lag (days)},ylabel={autocorrelation},xmin=0,xmax=250,ymin=0.6,ymax=1,title={\small spread level}]
\addplot[omDef,very thick] table[x=lag,y=levels]{figdata/methods/17-linear-time-series/acf.csv};
\nextgroupplot[xlabel={lag (days)},xmin=0,xmax=40,ymin=-0.06,ymax=0.06,title={\small daily change},
  scaled y ticks=false,yticklabel style={/pgf/number format/fixed,/pgf/number format/precision=2}]
\addplot[ycomb,omDef,thick,mark=*,mark size=1pt] table[x=lag,y=changes]{figdata/methods/17-linear-time-series/acf.csv};
\addplot[black,dashed,sharp plot] coordinates {(0,0.0178) (40,0.0178)};
\addplot[black,dashed,sharp plot] coordinates {(0,-0.0178) (40,-0.0178)};
\end{groupplot}
\end{tikzpicture}

\omcaption{Sample autocorrelations of the 2s10s Treasury spread, 1 June 1976 to 22 September 2026 (12\,574 days). Left: the level, lags 1 to
250, from 0.9987 down. Right: the daily change, lags 1 to 40, with the $\pm 2/\sqrt n$ band. Data: FRED series DGS10 and DGS2 (Board of
Governors of the Federal Reserve System, H.15).}\label{fig:qm:linear-time-series:acf}
\end{omfigure}

\section{Autoregressive and moving-average models}

\begin{definition}[Autoregressive, moving-average and ARMA processes]\label{def:qm:linear-time-series:arma}
With $(\varepsilon_t)$ \omterm{def:qm:linear-time-series:stationary}{white noise} of variance $\sigma^2$: an \emph{autoregressive process}\index{autoregressive process} of order $p$, AR($p$), satisfies
$X_t = c + \sum_{i=1}^p\phi_iX_{t-i} + \varepsilon_t$; a \emph{moving-average process}\index{moving-average process} of order $q$, MA($q$), is $X_t = \mu +
\varepsilon_t + \sum_{j=1}^q\vartheta_j\varepsilon_{t-j}$; an \emph{ARMA process}\index{ARMA process} combines them, $\phi(L)X_t = c + \vartheta(L)\varepsilon_t$ with
$\phi(z) = 1 - \sum_i\phi_iz^i$, $\vartheta(z) = 1 + \sum_j\vartheta_jz^j$ and $L$ the lag operator.
\end{definition}

\begin{proposition}[Stationarity and invertibility of ARMA]\label{prop:qm:linear-time-series:roots}
The ARMA equation has a unique weakly stationary causal solution $X_t = \mu + \sum_{j \ge 0}\psi_j\varepsilon_{t-j}$ if and only if every root of $\phi(z)$
lies outside the unit circle; it is invertible, $\varepsilon_t = \sum_{j \ge 0}\pi_jX_{t-j}$ up to a constant, if and only if every root of $\vartheta(z)$
does, the polynomials having no common root.
\end{proposition}

\begin{proof}
Factor $\phi(z) = \prod_i(1 - z/z_i)$. For $|z_i| > 1$, $(1 - L/z_i)^{-1} = \sum_kz_i^{-k}L^k$ is an absolutely summable filter, so $\psi(z) = \vartheta(z)/\phi(z)$
has summable coefficients and defines the solution; if some $|z_i| \le 1$ no summable causal inverse exists (for $|z_i| = 1$, none at all). The MA side
is the same argument applied to $\vartheta$.
\end{proof}

For the AR(1) $X_t = c + \phi X_{t-1} + \varepsilon_t$, the condition is $|\phi| < 1$, the autocorrelations are $\phi^h$, and a deviation from the mean
decays by half in $\ln(1/2)/\ln\phi$ periods, the \omterm{def:qm:stochastic-differential-equations:ou}{half-life} of chapter~4's \omterm{def:qm:stochastic-differential-equations:ou}{Ornstein--Uhlenbeck process} sampled daily. An AR($p$) has a \omterm{def:qm:linear-time-series:acf}{partial
autocorrelation function} that cuts off after lag $p$; an MA($q$) has an \omterm{def:qm:linear-time-series:acf}{autocorrelation function} that cuts off after lag $q$. The orders are
chosen by an \omterm{def:qm:linear-time-series:ic}{information criterion}.

\begin{definition}[Information criterion]\label{def:qm:linear-time-series:ic}
An \emph{information criterion}\index{information criterion} ranks models by fit penalised for size: Akaike's $\mathrm{AIC} = -2\ell + 2k$ and Schwarz's
$\mathrm{BIC} = -2\ell + k\ln n$, for maximised log-likelihood $\ell$ and $k$ parameters; the smallest wins.
\end{definition}

AIC aims at prediction and tends to choose large models; BIC is consistent for the true order when there is one. On the 2s10s daily changes AIC
chooses an AR(10) whose coefficients are all below 0.035 in absolute value: statistically detectable with 12\,574 days, economically nothing.

\section{The spectral view}

\begin{definition}[Spectral density, periodogram]\label{def:qm:linear-time-series:spectral}
The \emph{spectral density}\index{spectral density} of a weakly \omterm{def:qm:linear-time-series:stationary}{stationary process} with summable autocovariances is $f(\omega) = \frac1{2\pi}\sum_h\gamma(h)e^{-ih\omega}$,
$\omega \in [-\pi, \pi]$: the decomposition of its variance by frequency, $\gamma(0) = \int_{-\pi}^\pi f$. The \emph{periodogram}\index{periodogram} $I(\omega_j) =
\frac1{2\pi n}\bigl\lvert\sum_tX_te^{-it\omega_j}\bigr\rvert^2$ at the Fourier frequencies $\omega_j = 2\pi j/n$ is its sample counterpart.
\end{definition}

\omterm{def:qm:linear-time-series:stationary}{White noise} has a flat spectrum, $\sigma^2/2\pi$; an AR(1) with $\phi$ near one concentrates its variance at low frequencies, $f(\omega) =
\sigma^2/(2\pi|1 - \phi e^{-i\omega}|^2)$. The \omterm{def:qm:linear-time-series:spectral}{periodogram} is asymptotically unbiased but not consistent: each ordinate is approximately $f(\omega_j)$ times
an exponential variable, however long the sample, so it must be smoothed or averaged. Its low-frequency end is where persistence shows, and it is
the basis of the long-memory estimate of the last section. $f(0) = \frac1{2\pi}\sum_h\gamma(h)$ is, up to $2\pi$, the long-run variance of
chapter~11.

\section{Unit roots}

\begin{definition}[Unit root, spurious regression]\label{def:qm:linear-time-series:unitroot}
A process has a \emph{unit root}\index{unit root} if its autoregressive polynomial has a root at $z = 1$: its first difference is stationary but the level is
not, like a random walk. A \emph{spurious regression}\index{spurious regression} is a regression between independent processes with unit roots whose usual
$t$-statistics and $R^2$ indicate a relation that does not exist.
\end{definition}

Granger and Newbold (1974) showed the problem by simulation, and it is easy to repeat: regressing one random walk of 500 steps on another,
independent one, the slope's $t$-statistic exceeds 1.96 in absolute value in 89\% of 5\,000 trials, with a median $R^2$ of 0.17. The residual
inherits the \omterm{def:qm:linear-time-series:unitroot}{unit root}, the observations are not independent, and every formula that assumed they were is wrong. The same failure hits the
test of $\phi = 1$ itself.

\begin{definition}[Dickey--Fuller test]\label{def:qm:linear-time-series:df}
The \emph{Dickey--Fuller test}\index{Dickey--Fuller test} of a \omterm{def:qm:linear-time-series:unitroot}{unit root} regresses $\Delta X_t$ on $X_{t-1}$ (with a constant, and a trend if wanted) and
compares the $t$-statistic $\tau$ of the coefficient on $X_{t-1}$ with the Dickey--Fuller distribution rather than the normal; the augmented version
adds lagged differences $\Delta X_{t-k}$ to absorb short-run dependence (Said and Dickey, 1984).
\end{definition}

Under the \omterm{def:qm:linear-time-series:unitroot}{unit root} and with a constant, $\tau$ converges to $\bigl(\frac12(W(1)^2 - 1) - W(1)\int_0^1W\bigr)/\bigl(\int_0^1W^2 - (\int_0^1W)^2\bigr)^{1/2}$ for a
standard \omterm{def:qm:brownian-motion:bm}{Brownian motion} $W$, a functional that is neither normal nor centred at zero (\cref{fig:qm:linear-time-series:df}). Its quantiles have no closed
form; MacKinnon's (2010) response surfaces give the 1\%, 5\% and 10\% critical values $-3.43$, $-2.86$ and $-2.57$ for large samples. The chapter's own
simulation of 20\,000 random walks of 500 steps gives $-3.46$, $-2.89$ and $-2.57$, against MacKinnon's $-3.44$, $-2.87$ and $-2.57$ at that length.

\begin{omfigure}
\begin{tikzpicture}
\begin{axis}[width=10cm,height=5cm,xlabel={$t$-statistic of the coefficient on $X_{t-1}$},ylabel={density},
  tick label style={font=\small},label style={font=\small},xmin=-5,xmax=3,ymin=0,ymax=0.55,grid=major,grid style={gray!25},
  legend columns=2,legend style={font=\footnotesize,at={(0.5,-0.3)},anchor=north,draw=none,fill=none,
  /tikz/every even column/.append style={column sep=8pt}},table/col sep=comma]
\addplot[omDef,thick,const plot mark mid,fill=omDef!15] table[x=tau,y=density]{figdata/methods/17-linear-time-series/df.csv};
\addplot[omThm,very thick,dashed] table[x=tau,y=normal]{figdata/methods/17-linear-time-series/df.csv};
\addplot[black,thin,sharp plot] coordinates {(-2.86,0) (-2.86,0.55)};
\addplot[black,thin,dotted,sharp plot] coordinates {(-1.645,0) (-1.645,0.55)};
\legend{Dickey--Fuller (simulated),normal}
\end{axis}
\end{tikzpicture}

\omcaption{The law of the Dickey--Fuller $t$-statistic (regression with a constant) under a \omterm{def:qm:linear-time-series:unitroot}{unit root}, from 20\,000 simulated random walks of 500 steps,
against the standard normal. The 5\% critical values are $-2.86$ (solid line) and $-1.645$ (dotted): a test with the normal table rejects a true \omterm{def:qm:linear-time-series:unitroot}{unit
root} far too often. Data: the chapter's tutorial, seeded.}\label{fig:qm:linear-time-series:df}
\end{omfigure}

For the 2s10s spread, $\tau = -2.79$ with no lagged differences: not significant at 5\%, significant at 10\%. The augmented test depends on the lags:
$-2.74$ with five, $-3.13$ with ten, $-3.55$ with twenty; AIC, searching up to forty, chooses 34 and $-3.14$. On the sample since 1990 the Dickey--Fuller
statistic is $-2.01$. The honest reading is that the data cannot decide: the spread mean-reverts too slowly for fifty years of daily data to tell it
from a random walk.

\begin{proposition}[Kendall's bias of the AR(1) coefficient]\label{prop:qm:linear-time-series:kendall}
For a stationary AR(1) with an estimated mean, the least-squares coefficient satisfies $\E[\hat\phi] - \phi = -(1 + 3\phi)/n + O(n^{-2})$ (Kendall, 1954).
\end{proposition}

The bias is small in the coefficient and large in the \omterm{def:qm:stochastic-differential-equations:ou}{half-life}, because the \omterm{def:qm:stochastic-differential-equations:ou}{half-life} divides by $\ln\phi$, close to zero. The spread's estimate
0.99875, corrected by $(1 + 3 \times 0.99875)/12\,574 = 0.00032$, becomes 0.99906, and the \omterm{def:qm:stochastic-differential-equations:ou}{half-life} grows from 552 trading days (2.2 years) to 739 (2.9
years). The correction is an average: simulating 4\,000 histories of 12\,574 days from an AR(1) with $\phi = 0.99906$, the estimated \omterm{def:qm:stochastic-differential-equations:ou}{half-life} has median
576 days and a 10--90\% range of 352 to 969 days (\cref{fig:qm:linear-time-series:halflife}), and in 57\% of the histories the \omterm{def:qm:linear-time-series:df}{Dickey--Fuller test} does
not reject the \omterm{def:qm:linear-time-series:unitroot}{unit root} at 5\%, although there is none.

\begin{omfigure}
\begin{tikzpicture}
\begin{axis}[width=10cm,height=5cm,xlabel={estimated half-life (trading days)},ylabel={density ($\times 10^{-3}$)},
  tick label style={font=\small},label style={font=\small},xmin=0,xmax=2000,ymin=0,ymax=2.2,grid=major,grid style={gray!25},
  xtick={0,500,1000,1500,2000},xticklabel style={/pgf/number format/set thousands separator={}},
  legend style={font=\footnotesize,at={(0.98,0.96)},anchor=north east,fill=white,draw=none},table/col sep=comma]
\addplot[omDef,thick,const plot mark mid,fill=omDef!15] table[x=days,y=density]{figdata/methods/17-linear-time-series/halflife.csv};
\addplot[black,dashed,thick,sharp plot] coordinates {(739,0) (739,2.2)};
\addplot[omThm,very thick,sharp plot] coordinates {(552,0) (552,2.2)};
\legend{simulated estimates,true half-life (739),2s10s estimate (552)}
\end{axis}
\end{tikzpicture}

\omcaption{Sampling distribution of the AR(1) \omterm{def:qm:stochastic-differential-equations:ou}{half-life} estimated from 12\,574 daily observations of an AR(1) whose true \omterm{def:qm:stochastic-differential-equations:ou}{half-life} is 739 days, over
4\,000 simulated histories (estimates above 2\,000 days piled in the last bin): median 576, 10--90\% range 352 to 969. The 2s10s estimate of 552 days is a
typical draw. Data: the chapter's tutorial, seeded.}\label{fig:qm:linear-time-series:halflife}
\end{omfigure}

\section{Long memory}

\begin{definition}[Long memory, fractional differencing, Hurst exponent]\label{def:qm:linear-time-series:longmem}
A \omterm{def:qm:linear-time-series:stationary}{stationary process} has \emph{long memory}\index{long memory} if its autocorrelations decay like a power, $\rho(h) \sim ch^{2d-1}$ with $0 < d < \frac12$, so
that they are not summable. \emph{Fractional differencing}\index{fractional differencing} applies $(1 - L)^d = \sum_{j \ge 0}w_jL^j$ with $w_0 = 1$, $w_j = w_{j-1}(j - 1 -
d)/j$; the ARFIMA$(0, d, 0)$ process is $(1 - L)^dX_t = \varepsilon_t$ (Granger and Joyeux, 1980; Hosking, 1981). The \emph{Hurst exponent}\index{Hurst exponent} is
$H = d + \frac12$, measured by the growth of ranges or variances of partial sums like $n^H$ (Hurst, 1951).
\end{definition}

\begin{definition}[Fractional Brownian motion]\label{def:qm:linear-time-series:fbm}
\emph{Fractional Brownian motion}\index{fractional Brownian motion} with \omterm{def:qm:linear-time-series:longmem}{Hurst exponent} $H \in (0, 1)$ is the centred \omterm{def:qm:brownian-motion:bm}{Gaussian process} with $B_H(0) = 0$ and
$\E[B_H(t)B_H(s)] = \frac12(t^{2H} + s^{2H} - |t - s|^{2H})$ (Mandelbrot and Van Ness, 1968). For $H = \frac12$ it is \omterm{def:qm:brownian-motion:bm}{Brownian motion}; for $H > \frac12$ its increments
are positively correlated with power-law decay, for $H < \frac12$ negatively.
\end{definition}

For $H \ne \frac12$ \omterm{def:qm:linear-time-series:fbm}{fractional Brownian motion} is not a \omterm{def:qm:jump-processes:jd}{semimartingale}, so the Itô calculus of chapter~3 does not apply to it, and a price that followed it
would admit arbitrage; it is a model for volatility and for slowly mean-reverting quantities, not for prices. The difference between long and short
memory is in the tail of the correlogram (\cref{fig:qm:linear-time-series:longmem}): an ARFIMA with $d = 0.3$ and an AR(1) with the same first
autocorrelation, 0.41, differ by a factor of a thousand at lag 10. On 20\,000 simulated observations the rescaled-range estimate of $H$ is 0.78 and the
\omterm{def:qm:linear-time-series:spectral}{log-periodogram} (Geweke--Porter-Hudak) estimate 0.71, for a true 0.8; on \omterm{def:qm:linear-time-series:stationary}{white noise} they give 0.53 and 0.46, the rescaled range being biased upward in short
blocks.

\begin{omfigure}
\begin{tikzpicture}
\begin{loglogaxis}[width=10cm,height=5cm,xlabel={lag},ylabel={autocorrelation},
  tick label style={font=\small},label style={font=\small},xmin=1,xmax=200,ymin=1e-4,ymax=1,grid=major,grid style={gray!25},
  legend columns=3,legend style={font=\footnotesize,at={(0.5,-0.3)},anchor=north,draw=none,fill=none,
  /tikz/every even column/.append style={column sep=8pt}},table/col sep=comma]
\addplot[only marks,mark=*,mark size=0.9pt,omDef] table[x=lag,y=arfima]{figdata/methods/17-linear-time-series/longmem.csv};
\addplot[omDef,very thick] table[x=lag,y=theory]{figdata/methods/17-linear-time-series/longmem.csv};
\addplot[omThm,very thick,dashed] table[x=lag,y=ar1]{figdata/methods/17-linear-time-series/longmem.csv};
\legend{{ARFIMA, $d = 0.3$, sample},{ARFIMA, theory},{AR(1), same $\rho(1)$}}
\end{loglogaxis}
\end{tikzpicture}

\omcaption{Autocorrelations of a long-memory process (ARFIMA with $d = 0.3$: 20\,000 simulated observations and the exact $\rho(h) = \prod_{i \le h}(i - 1 + d)/(i -
d)$, a straight line of slope $2d - 1$ on these axes) and of an AR(1) with the same lag-one autocorrelation, which decays geometrically. Data: the chapter's
tutorial, seeded.}\label{fig:qm:linear-time-series:longmem}
\end{omfigure}

The spread's daily changes give a \omterm{def:qm:linear-time-series:spectral}{log-periodogram} estimate $H = 0.34$ from the lowest $\sqrt n = 112$ frequencies, that is $d = -0.16$ with a \omterm{def:qm:estimation:estimator}{standard error}
of $\pi/\sqrt{24 \times 112} = 0.06$: the changes are over-differenced, and the level behaves like a fractionally integrated process with $d \approx 0.84$, mean
reverting in the long run but with hyperbolic rather than geometric decay. That is one reason AR(1) half-lives of the spread are unstable across samples:
the AR(1) is the wrong shape of memory.

\section{Tutorial: the half-life of a spread}\label{tut:qm:linear-time-series:halflife}

\begin{tutorial}
\textbf{Goal.} Estimate the 2s10s spread's persistence three ways and decide what the data can say about a \omterm{def:qm:linear-time-series:unitroot}{unit root}.
\textbf{End state:} \cref{fig:qm:linear-time-series:acf,fig:qm:linear-time-series:df,fig:qm:linear-time-series:halflife}; the half-lives 552 and 739 days and the
Dickey--Fuller statistics.
\begin{tutsteps}
\item \textbf{The AR(1) with its \omterm{def:qm:stochastic-differential-equations:ou}{half-life} and Kendall's correction.}
\omcode{code/firm/tsa/firm_tsa.py}{102}{118}{AR(1) by least squares, half-life and bias correction.}\label{lst:qm:linear-time-series:ar1}
\item \textbf{The augmented Dickey--Fuller regression and MacKinnon's critical values.}
\omcode{code/firm/tsa/firm_tsa.py}{125}{158}{The augmented Dickey--Fuller test.}\label{lst:qm:linear-time-series:adf}
\item \textbf{Run} \texttt{problem()}, \texttt{half\_life\_mc(0.99906, 12574)}, \texttt{spurious()} and \texttt{long\_memory()} in \texttt{qm\_tsa.py}, then
\texttt{fig\_tsa.py}.
\end{tutsteps}
\textbf{What to change next.} Fit an ARFIMA to the spread by the \omterm{def:qm:linear-time-series:spectral}{log-periodogram} $d$ and compare its forecast of the spread's decay with the AR(1)'s; rerun the
tests on weekly data and see which conclusions survive.
\end{tutorial}

\section{Build: the time-series toolkit}\label{bld:qm:linear-time-series:tsa}

\begin{build}
\bfield{Purpose} Every persistence the miniature firm relies on (a spread's \omterm{def:qm:stochastic-differential-equations:ou}{mean reversion}, a signal's decay, a volatility's memory) is measured and tested
here before a strategy is built on it.
\bfield{Interface} \texttt{acf}, \texttt{pacf}, \texttt{band}; \texttt{yule\_walker(x, p)}; \texttt{arma\_css(x, p, q)}; \texttt{ar1(x)}; \texttt{adf(x, lags,
trend, max\_lags)}; \texttt{mackinnon\_crit(trend, n)}; \texttt{periodogram(x)}; \texttt{frac\_diff(x, d, k)}; \texttt{hurst\_rs(x)}, \texttt{hurst\_gph(x)}.
\bfield{Rules} Unit-root tests report their lag choice and are rerun over a range of lags; half-lives are reported with the bias correction and a simulated
interval; no normal table for a coefficient near one.
\bfield{Acceptance tests} \texttt{code/firm/tsa/tests/}: autocorrelations and partial autocorrelations of an AR(1); Yule--Walker for an AR(2); ARMA(1, 1) by
conditional least squares; the \omterm{def:qm:linear-time-series:df}{Dickey--Fuller test} has its nominal size on random walks and rejects a stationary AR(1); MacKinnon's formula; a flat \omterm{def:qm:linear-time-series:spectral}{periodogram} for
\omterm{def:qm:linear-time-series:stationary}{white noise}; \omterm{def:qm:linear-time-series:longmem}{fractional differencing} with $d = 1$ is differencing; Hurst estimates of \omterm{def:qm:linear-time-series:stationary}{white noise} and of an ARFIMA.
\bfield{Stretch} The KPSS test, whose null is stationarity; exact Gaussian likelihood for ARMA by the Kalman filter of chapter~19; the local Whittle \omterm{def:qm:estimation:estimator}{estimator} of $d$.
\end{build}

\begin{omsources}
\item H.~Wold, \emph{A Study in the Analysis of Stationary Time Series}, Almqvist and Wiksell, 1938.
\item D.~A.~Dickey and W.~A.~Fuller, ``Distribution of the estimators for autoregressive time series with a unit root'', \emph{Journal of the American Statistical
Association} 74, 1979; S.~E.~Said and D.~A.~Dickey, \emph{Biometrika} 71, 1984.
\item J.~G.~MacKinnon, ``Critical values for cointegration tests'', Queen's Economics Department Working Paper 1227, 2010.
\item C.~W.~J.~Granger and P.~Newbold, ``Spurious regressions in econometrics'', \emph{Journal of Econometrics} 2, 1974.
\item M.~G.~Kendall, ``Note on bias in the estimation of autocorrelation'', \emph{Biometrika} 41, 1954.
\item H.~E.~Hurst, ``Long-term storage capacity of reservoirs'', \emph{Transactions of the American Society of Civil Engineers} 116, 1951.
\item B.~B.~Mandelbrot and J.~W.~Van Ness, ``Fractional Brownian motions, fractional noises and applications'', \emph{SIAM Review} 10, 1968.
\item C.~W.~J.~Granger and R.~Joyeux, \emph{Journal of Time Series Analysis} 1, 1980; J.~R.~M.~Hosking, ``Fractional differencing'', \emph{Biometrika} 68, 1981;
J.~Geweke and S.~Porter-Hudak, \emph{Journal of Time Series Analysis} 4, 1983.
\item Board of Governors of the Federal Reserve System, H.15 Selected Interest Rates, via FRED (series DGS2 and DGS10), accessed 24 September 2026.
\end{omsources}

\section{Exercises}

\begin{exercise}[$\star$]\label{exo:qm:linear-time-series:1}
An AR(1) has coefficient 0.98 on daily data. What is the \omterm{def:qm:stochastic-differential-equations:ou}{half-life} in days, and what is the autocorrelation at lag 20?
\end{exercise}

\begin{exercise}[$\star$]\label{exo:qm:linear-time-series:2}
Is the MA(1) $X_t = \varepsilon_t + 2\varepsilon_{t-1}$ invertible? Find the invertible MA(1) with the same autocorrelations.
\end{exercise}

\begin{exercise}[$\star$]\label{exo:qm:linear-time-series:3}
Write the first four fractional-differencing weights for $d = 0.3$.
\end{exercise}

\begin{exercise}[$\star\star$]\label{exo:qm:linear-time-series:4}
For an AR(1) with $\phi = 0.99$ estimated from 1\,000 daily observations, what is Kendall's bias, and how much does it change the \omterm{def:qm:stochastic-differential-equations:ou}{half-life}?
\end{exercise}

\begin{exercise}[$\star\star$]\label{exo:qm:linear-time-series:5}
Show that the autocorrelations of an AR(2) satisfy $\rho(h) = \phi_1\rho(h - 1) + \phi_2\rho(h - 2)$, and compute $\rho(1)$ and $\rho(2)$ for $\phi_1 = 0.5$, $\phi_2 =
-0.3$.
\end{exercise}

\begin{exercise}[$\star\star$]\label{exo:qm:linear-time-series:6}
Using MacKinnon's coefficients $(-2.86154, -2.8903, -4.234, -40.040)$, compute the 5\% Dickey--Fuller critical value (with constant) for $T = 100$ and $T = 12\,574$.
\end{exercise}

\begin{exercise}[$\star\star\star$]\label{exo:qm:linear-time-series:7}
\emph{Coding.} Regress one independent random walk of 500 steps on another 5\,000 times. How often is $|t| > 1.96$? What happens if you regress the differences
instead?
\end{exercise}

\begin{exercise}[$\star\star\star$]\label{exo:qm:linear-time-series:8}
\emph{Find the flaw.} ``We regressed the daily change of the spread on the lagged level; the coefficient is $-0.00125$ with $t = -2.79$, significant at 1\%, so the
spread mean-reverts and we will trade its reversion with a \omterm{def:qm:stochastic-differential-equations:ou}{half-life} of 552 days.''
\end{exercise}

\section{Problem: The Half-Life of a Spread}

\begin{problem}[{Weekend problem --- how fast does the 2s10s spread mean-revert?}]\label{pb:qm:linear-time-series:1}
The data are the daily ten-year and two-year constant-maturity Treasury yields (FRED DGS10 and DGS2) from 1 June 1976 to 22 September 2026, on the 12\,574
days with both; the spread is their difference in basis points.

\medskip\noindent\textbf{Part I --- The series.}
\begin{enumerate}
\item What are the spread's mean, range and last value?
\item What is its first autocorrelation, and what does its correlogram look like?
\item What are the standard deviation and first autocorrelations of its daily changes?
\item Which AR order does AIC choose for the changes, and does it matter economically?
\item What is the Wold representation of the changes, approximately?
\end{enumerate}

\medskip\noindent\textbf{Part II --- The AR(1).}
\begin{enumerate}[resume]
\item What are the AR(1) coefficient, its \omterm{def:qm:estimation:estimator}{standard error}, and the \omterm{def:qm:stochastic-differential-equations:ou}{half-life}?
\item What $t$-statistic against one, and what $p$-value by the normal table?
\item What is Kendall's correction, and the corrected \omterm{def:qm:stochastic-differential-equations:ou}{half-life}?
\item What do the Dickey--Fuller and augmented tests say, with 0, 5, 10 and 20 lags?
\item What does the sample since 1990 give?
\end{enumerate}

\medskip\noindent\textbf{Part III --- What the data can tell.}
\begin{enumerate}[resume]
\item If the true \omterm{def:qm:stochastic-differential-equations:ou}{half-life} were 739 days, how would 12\,574-day estimates be distributed?
\item How often would the \omterm{def:qm:linear-time-series:df}{Dickey--Fuller test} then fail to reject the \omterm{def:qm:linear-time-series:unitroot}{unit root}?
\item What does the \omterm{def:qm:linear-time-series:spectral}{log-periodogram} say about the changes, and about the level?
\item Why might an AR(1) \omterm{def:qm:stochastic-differential-equations:ou}{half-life} be unstable across samples?
\item What would a steepener trader want to know that the \omterm{def:qm:stochastic-differential-equations:ou}{half-life} does not say?
\end{enumerate}

\medskip\noindent\textbf{Part IV --- Judgement.}
\begin{enumerate}[resume]
\item Is the spread stationary?
\item Which \omterm{def:qm:stochastic-differential-equations:ou}{half-life} would you use to size a mean-reversion trade, and with what caveat?
\item Why is the normal table's $p$-value of 0.003 misleading?
\item State the \emph{named result}: the AR(1) \omterm{def:qm:stochastic-differential-equations:ou}{half-life}, its bias-corrected value, and whether a \omterm{def:qm:linear-time-series:unitroot}{unit root} can be rejected.
\item In one sentence: what does a unit-root test tell a trader?
\end{enumerate}
\end{problem}

\section{Interview questions}

\begin{interviewq}[$\star$ \iqroles{researcher, trader}]\label{iq:qm:linear-time-series:1}
What is the \omterm{def:qm:stochastic-differential-equations:ou}{half-life} of an AR(1) with coefficient $\phi$, and why is it hard to estimate when $\phi$ is near one?
\end{interviewq}

\begin{interviewq}[$\star\star$ \iqroles{researcher}]\label{iq:qm:linear-time-series:2}
Why can't you use the usual $t$-table to test $\phi = 1$?
\end{interviewq}

\begin{interviewq}[$\star\star$ \iqroles{researcher, mle}]\label{iq:qm:linear-time-series:3}
You regress one price series on another and get $R^2 = 0.6$ and $t = 15$. What do you check?
\end{interviewq}

\begin{interviewq}[$\star\star$ \iqroles{researcher}]\label{iq:qm:linear-time-series:4}
How do you tell an AR process from an MA process from their correlograms?
\end{interviewq}

\begin{interviewq}[$\star\star$ \iqroles{researcher, mle}]\label{iq:qm:linear-time-series:5}
What is \omterm{def:qm:linear-time-series:longmem}{long memory}, and where do you see it in markets?
\end{interviewq}

\begin{interviewq}[$\star\star\star$ \iqroles{researcher}]\label{iq:qm:linear-time-series:6}
Why is the least-squares AR(1) coefficient biased downward in small samples, and how large is the bias?
\end{interviewq}
